\section{An actual joint-cutoff Gram remainder}
\label{sec:joint-cutoff-counterpath}

The source's two divergent cutoffs do not specify a relative rate.
The following construction keeps the actual global PW projection,
the actual Gram-normalized transport, and the repaired finite-place
local quotient as three explicitly defined operators.  It gives a
nonzero remainder along an admissible joint path.  The argument
does not evaluate the remainder by subtracting two zero sums.
It first constructs a nonzero trace on the actual quotient and
then uses trace-ideal convergence of the actual operators.

Throughout this section take the rational trivial sector,
$\alpha>0$, $s\ge0$, and
\[
 \begin{aligned}
 H&=L^2(\mathbb R,w(x)\,dx),&
 w(x)&=e^{2\alpha|x|}(1+x^2)^{s/2},\\
 J&=\overline{I_{1/2}}^{\,H},&
 Z&=J^\perp,&
 Q_Y^{\rm glob}&=P_{J_Y}.
 \end{aligned}
 \tag{MA-JP.1}
\]
The notation $I_{1/2}$ is the actual centered smooth theta ideal,
and $J_Y$ is the actual finite Fourier-band theta range of
MA-GP-001--008.  The quotient trace from MA-WP is denoted by
\[
 \tau(g)=\operatorname{Tr}_{H/J}C_g^{H/J}
        =\operatorname{Tr}_H(C_gP_Z),
 \qquad
 G(z)=\int_{\mathbb R}g(x)e^{zx}\,dx .
 \tag{MA-JP.2}
\]
The second equality is the actual ambient lift and trace cyclicity
of MA-WP/MA-LIFT, not a definition through a scalar divisor.

\paragraph{MA-JP-001 (the exact source rate conditions).}
Masullo, Definition~9.3.3.1, printed p.~427, specifies a
multiplicatively symmetric cutoff supported on $[1/X,X]$,
equal to one on $[1/X+\delta,X-\delta]$.
At the end of the local calculation, printed p.~519,
(12.4.1.25)--(12.4.1.28), the limits stated are
\[
       X,Y\longrightarrow\infty,\qquad
       \delta X\longrightarrow0,\qquad
       \delta Y\longrightarrow0 .
 \tag{MA-JP.3}
\]
Here $X=\lambda$, $Y=\lambda'$, and $\delta$ is the source's
boundary-smoothing parameter.  It is not the lower spatial
endpoint.  The footnote on p.~514 explicitly proposes the
annular replacement $[1/\lambda'',\lambda]$ with independent
$\lambda''$, and requests the convergence study of its volume
integral.  Neither that footnote nor the conditions on p.~519
prescribe $Y/\lambda''\to0$.  For the symmetric cutoff of
Definition~9.3.3.1, $\lambda''=X$, so the parameter of
MA-LA-014 is
\[
                     \epsilon Y=Y/X.
 \tag{MA-JP.4}
\]
The proof of Theorem~12.5.1.1, p.~520, invokes the local
calculation for the global geometric kernel but supplies no
additional relation between these two cutoffs or a quantitative
Gram estimate.  This is a precise reading of the cited passages,
not a claim that no possible additional scaling could work.

\paragraph{MA-JP-002 (a compact witness with both polar moments removed).}
Suppose the retained quotient divisor of MA-WP is nonempty:
$r_{\rho_0}>0$ for some nontrivial completed zero $\rho_0$.
Then there is a test $g\in C_c^\infty(\mathbb R)$ such that
\[
        g(0)=0,\qquad G(1/2)=G(-1/2)=0,\qquad \tau(g)=1 .
 \tag{MA-JP.5}
\]
The normalization of the last value is a harmless complex
scalar normalization of the test.  In particular its polar
term $P(g)=G(1/2)+G(-1/2)$ is zero.

\begin{proof}
Put $\lambda_0=\rho_0-1/2$ and
$q(z)=z^2-1/4$.  Because $\rho_0$ is a nontrivial completed
zero, $\lambda_0\ne\pm1/2$.
The actual full-jet interpolation of MT-M18 supplies
$f\in\mathcal E$ whose centered Mellin value at $\rho_0$
is $1/(r_{\rho_0}q(\lambda_0))$ and whose jets at all
other zeros are zero.  The finitely many jets in the
chosen block can be prescribed by the finite-dimensional
Hermite classes; no assertion about individual zero
separation is needed.
The differential operator
\[
                     \mathcal D=\frac{d^2}{dx^2}-\frac14
 \tag{MA-JP.6}
\]
multiplies the Mellin transform by $q(z)$.
Consequently $f_1=\mathcal D f$ has both polar moments
zero and $\tau(f_1)=1$.  Here the finite-block trace
$r_\rho G(\rho-1/2)$ is used only to choose a nonzero
vector in the already constructed actual trace functional.
Continuity of that functional on $\mathcal E$ follows
from the rapid block estimates of MA-WP, or directly
from the zero count and rapid Mellin decay.

There is $v\in I_{1/2}$ with both polar moments zero and
$v(0)=1$.  To see this, choose any nonzero
$\theta\in I_{1/2}$.  Then $\mathcal D\theta\ne0$:
the only smooth solutions of $\mathcal D\theta=0$
are linear combinations of $e^{x/2}$ and $e^{-x/2}$,
and neither nonzero combination belongs to $\mathcal E$.
The ideal is stable under differentiation and translation,
as is immediate from its full Mellin-jet description.
Translate $\mathcal D\theta$ to put a nonzero value at
zero, and divide by that value.
This produces $v$.  In particular $\tau(v)=0$.
Thus
\[
                        f_2=f_1-f_1(0)v
 \tag{MA-JP.7}
\]
still has the two zero polar moments and trace one,
and now vanishes at zero.

Choose $\kappa\in C_c^\infty(\mathbb R)$ equal to one
near zero and put $f_N=\kappa(x/N)f_2(x)$.
Then $f_N\to f_2$ in $\mathcal E$ and $f_N(0)=0$.
To correct both small polar errors exactly, choose a
nonnegative smooth bump $b$ supported near zero and
two distinct positive numbers $x_1,x_2$ so that
$b_j(x)=b(x-x_j)$ are supported away from zero.
The matrix
\[
 \begin{pmatrix}
 \mathcal Mb_1(1/2)&\mathcal Mb_2(1/2)\\
 \mathcal Mb_1(-1/2)&\mathcal Mb_2(-1/2)
 \end{pmatrix}
 =
 \begin{pmatrix}
 e^{x_1/2}\mathcal Mb(1/2)&e^{x_2/2}\mathcal Mb(1/2)\\
 e^{-x_1/2}\mathcal Mb(-1/2)&e^{-x_2/2}\mathcal Mb(-1/2)
 \end{pmatrix}
 \tag{MA-JP.8}
\]
is invertible.  Its determinant is a nonzero constant
times $2\sinh((x_1-x_2)/2)$.
Let $(c_{1,N},c_{2,N})$ solve this fixed system with
right side $(\mathcal Mf_N(1/2),\mathcal Mf_N(-1/2))$.
These coefficients tend to zero.  Hence
$u_N=f_N-c_{1,N}b_1-c_{2,N}b_2$ converges to $f_2$
in $\mathcal E$, vanishes at zero, and has both polar
moments exactly zero.  For all sufficiently large $N$,
$\tau(u_N)\ne0$ by continuity.  Set
$g=u_N/\tau(u_N)$.  This proves (MA-JP.5) with an
actual compact smooth test.
\end{proof}

\paragraph{The classical zeta instance and its exact input.}
For every $\alpha>0$ and every $s\ge0$, the hypothesis of
MA-JP-002 is satisfied for zeta by the classical theorem
of Hardy (1914), which gives infinitely many zeros on
$\operatorname{Re}s=1/2$.  Each such zero lies strictly
inside the corrected strip and retains its full multiplicity.
Only existence of one such zero is needed.
This invocation is unconditional arithmetic information
about a zero's location; it is not an absence of all
zero-location input.  It assumes neither the Riemann
hypothesis nor the theorem $\mathsf A$ nor $M$.
The primary statement is Hardy's opening paragraph,
printed p.~1012, with its argument on pp.~1012--1014.

\subsection{Two local comparisons with their Hilbert metrics stated}

Write $h(t)=g(\log t)$ and let $S$ be any fixed finite set
of rational primes, together with the real place.  It may
be chosen to include every prime power norm in the support
of $h$, as required when the source passes to p.~520.
On $\mathscr H_e=L^2_{\rm even}(\mathbb R,dt)$ use the
actual local operators of MA-LA-011--014,
\[
 \begin{aligned}
 (A_hu)(t)&=\int_0^\infty h(r)r^{-1/2}u(t/r)\,\frac{dr}{r},\\
 Q_Y&=\mathcal F_e^{-1}M_{\mathbf1_{[-Y,Y]}}\mathcal F_e,&
 D_S&=\prod_{p\in S_f}D_p .
 \end{aligned}
 \tag{MA-JP.9}
\]
The finite-place product $D_S$ is the actual unitary
quotient Fourier factor, not an assigned prime scalar.
All factors in this paragraph have fixed $S$.
The logarithmic unitary
\[
 (\Gamma_\infty u)(x)=\sqrt2\,e^{x/2}u(e^x)
 \tag{MA-JP.10}
\]
sends $A_h$ to $C_g$ on unweighted $L^2(\mathbb R,dx)$.
The global weighted unitary of MA-GP-007 is instead
\[
 (\mathcal U_wu)(t)
       =\sqrt{\frac{2t}{w(\log t)}}u(t),\qquad t>0 .
 \tag{MA-JP.11}
\]
It maps $\mathscr H_e$ to the radial form of $H$.
For clarity, transporting global convolution by this
second unitary gives the distinct operator
\[
 \begin{split}
 (\mathcal U_w^{-1}C_g\mathcal U_wu)(t)
 =\int_0^\infty h(r)r^{-1/2}
       \sqrt{\frac{w(\log t)}{w(\log t-\log r)}}
                  u(t/r)\,\frac{dr}{r}.
 \end{split}
 \tag{MA-JP.12}
\]
Thus the actual local convolution and the same-norm
additive comparison must not be silently equated.
Both comparisons below are carried out explicitly.

Let $E_Y:\mathcal P_Y^0\to\mathscr H_e$ be the ordinary
band inclusion of MA-GP-008 and set
\[
 W_Y^S=\mathcal U_wD_S\mathcal F_eE_Y,\qquad
 \Pi_Y^{S,0}=W_Y^S(W_Y^S)^* .
 \tag{MA-JP.13}
\]
This is an orthogonal projection in the same weighted
space $H$ as $Q_Y^{\rm glob}$.
For $S_f=\varnothing$ it is exactly the zero-integral
additive candidate in (MA-GP.73).
It is a precisely specified comparison, not an
operator identity asserted to be printed in the source.

The finite set $S$ may also be retained in the \emph{global}
PW family itself.  Let $\widehat J_{S,Y}$ be the $H$-closure
of the actual global theta images of polar-free
$S$-adelic Schwartz functions whose additive Fourier
support has product norm at most $Y$, extended by
$\mathbf1_{\mathbb Z_p}$ outside $S$ and projected to
the trivial compact type.  This is the rational invariant
PW level of Definition~7.3.2.3, with no additional
finite transition set.  Write
$\widehat Q_{S,Y}=P_{\widehat J_{S,Y}}$.
There are literal inclusions
\[
         J_Y\subseteq\widehat J_{S,Y}\subseteq J,
       \qquad \widehat Q_{S,Y}\longrightarrow P_G
                          \quad\hbox{strongly}.
 \tag{MA-JP.13a}
\]
For the first inclusion use the actual vector
$f_\infty\otimes\bigotimes_{p\in S_f}\mathbf1_{\mathbb Z_p}$.
Its finite Fourier factors are again
$\mathbf1_{\mathbb Z_p}$, so their norms are at most one,
and an archimedean band $Y$ meets the stated product
support condition.  After extending by the outside
integer balls, the global rational summation is exactly
the same $2\sqrt t\sum_{n\ge1}f_\infty(nt)$ as before.
Both global polar conditions follow from those of
$f_\infty$.  The second inclusion holds because these
are actual global polar-free theta images in the
centered smooth ideal, using its Meyer/MT-M24
identification.  The levels increase, and the union
of the smaller $J_Y$ is dense in $J$ by MA-GP-001.
The orthogonal projection argument of that theorem
proves the last assertion.  This proves only the
inclusions and strong limit needed here; it does not
identify the mixed-place finite range with the
one-place Hardy model or its frequency Gram.

\paragraph{MA-JP-003 (a nonzero actual Gram remainder along a source-admissible path).}
For $g$ as in (MA-JP.5) and fixed $S$, there are
$X_n=n+2$, integers $Y_n\ge X_n^2$, and positive
$\delta_n\to0$ with
\[
 \delta_nX_n\longrightarrow0,\qquad
 \delta_nY_n\longrightarrow0,\qquad Y_n/X_n\longrightarrow\infty,
 \tag{MA-JP.14}
\]
and smooth, multiplicatively symmetric spatial cutoffs
$\rho_n$ of Definition~9.3.3.1, of level $(X_n,\delta_n)$,
and even smooth frequency cutoffs $b_n$ at $Y_n$, such that
the following \emph{actual} traces have the limits
\[
 \begin{aligned}
 \operatorname{Tr}_H(M_{\rho_n}C_gQ_{Y_n}^{\rm glob})
       &\longrightarrow-1,\\
 \operatorname{Tr}_H(M_{\rho_n}C_g\widehat Q_{S,Y_n})
       &\longrightarrow-1,\\
 \operatorname{Tr}_H(M_{\rho_n}C_g\Pi_{Y_n}^{S,0})
       &\longrightarrow0,\\
 \operatorname{Tr}_{\mathscr H_e}
   \bigl(M_{\rho_n(|t|)}A_hD_S
           \mathcal F_e^{-1}M_{b_n}\mathcal F_eD_S^*\bigr)
       &\longrightarrow0 .
 \end{aligned}
 \tag{MA-JP.15}
\]
In the first three lines, $M_{\rho_n}$ means multiplication
by $\rho_n(e^x)$ in logarithmic coordinates.
The frequency cutoff satisfies $0\le b_n\le1$, is
supported on $[-Y_n,Y_n]$, and is one on
$[-Y_n+\delta_n,Y_n-\delta_n]$.
The last line is the actual repaired finite-$S$ local
quotient operator, with its nuclear section lift.
Its coordinate realization, measure, and finite-place
factors are precisely those of MA-LA-014.
Because $P(g)=0$, subtracting the polar term does not
alter any of these limits.

In particular each global trace minus either local
comparison trace tends to $-1$.  The one-place global
trace minus the same-norm additive trace also has the
literal Gram-resolvent expression of (MA-GP.74) with
$W_{Y_n}$ replaced by the stated isometry $W_{Y_n}^S$.

\begin{proof}
\emph{Compact-position trace ideals, including sharp cutoffs.}
For $X>1$ let
$a_X(t)=\mathbf1_{[1/X,X]}(t)$ and
$\chi_X(x)=a_X(e^x)$.
The operators
\[
 A_X=M_{\chi_X}C_g\quad\hbox{on }H,\qquad
 B_X=M_{a_X(|t|)}A_h\quad\hbox{on }\mathscr H_e
 \tag{MA-JP.16}
\]
are trace class, and both have trace zero.
Indeed insert a fixed smooth cutoff equal to one
on the compact output interval, and a second smooth
cutoff on its enlargement by $\operatorname{supp}g$.
The resulting convolution kernel is smooth on a compact
rectangle and hence trace class by its absolutely
summable Fourier expansion.  Multiplication by the
bounded characteristic function preserves trace class.
On the weighted space, the factors $w^{1/2}$ and
$w^{-1/2}$ are bounded on the two coordinate
projections.  They need no smoothness at the cusp.
They cancel on the diagonal, as may also be checked
by cyclicity after these compact cutoffs are inserted.
Smooth bounded approximations of $\chi_X$ converge
strongly; multiplication of the fixed trace-class
compact-kernel operator then converges in trace norm.
Its diagonal traces converge to
$g(0)\int\chi_X(x)\,dx=0$.
Formula (MA-JP.10) proves the identical assertion for
$B_X$.  This argument also proves trace-norm convergence
when either sharp spatial cutoff is smoothed at fixed $X$.

As $X\to\infty$, multiplication by $\chi_X$ tends
strongly to the identity, with norm one.  The actual
ambient lift $C_gP_Z$ is trace class by MA-WP.
Consequently
\[
 \begin{split}
 \operatorname{Tr}_H(A_XP_G)
 &=\operatorname{Tr}_H A_X-
            \operatorname{Tr}_H(M_{\chi_X}C_gP_Z)\\
 &\longrightarrow-\operatorname{Tr}_H(C_gP_Z)=-1.
 \end{split}
 \tag{MA-JP.17}
\]
This is a limit of traces of actual trace-class
operators.  No spectral scalar sum has been
substituted for either cutoff operator.

\emph{The fixed-position convergence statements.}
For every fixed $X$, MA-GP-001 and finite-rank
trace-ideal approximation give
\[
            \|A_X(Q_Y^{\rm glob}-P_G)\|_1\longrightarrow0 .
 \tag{MA-JP.18}
\]
The same argument applies to the sharp $A_X$ just
proved trace class, since the PW projections are
self-adjoint and converge strongly.
It also gives
$\|A_X(\widehat Q_{S,Y}-P_G)\|_1\to0$
by (MA-JP.13a).

Also $\Pi_Y^{S,0}\to1$ strongly on $H$.
On the additive frequency space the band projection
onto zero-integral even functions is
\[
 P_Y^0=M_{\mathbf1_{[-Y,Y]}}-e_Y\otimes e_Y,\qquad
              e_Y=(2Y)^{-1/2}\mathbf1_{[-Y,Y]} .
 \tag{MA-JP.19}
\]
The first term converges strongly to one.
The unit vectors $e_Y$ converge weakly to zero:
this follows first for compactly supported $L^2$
functions by their $L^1$ bound, and then by density.
Thus their rank-one projections converge strongly
to zero.  Conjugation by the fixed unitaries in
(MA-JP.13) proves the claim.
Finally $D_SQ_YD_S^*\to1$ strongly on $\mathscr H_e$.
All these operators are self-adjoint contractions.
Therefore
\[
 \begin{aligned}
 \|A_X(\Pi_Y^{S,0}-1)\|_1&\longrightarrow0,\\
 \|B_X(D_SQ_YD_S^*-1)\|_1&\longrightarrow0.
 \end{aligned}
 \tag{MA-JP.20}
\]
The elementary trace-ideal fact used here follows
for a rank-one $u\otimes\ell$ from strong convergence
of the adjoints acting on $\ell$, and then for every
trace-class operator by finite-rank approximation.

\emph{Selecting one path before smoothing.}
For $X_n=n+2$, choose an integer $Y_n\ge X_n^2$
so large that each norm in (MA-JP.18) and
(MA-JP.20), and the corresponding norm for
$\widehat Q_{S,Y}$, at $X=X_n$, is at most $1/n$.
All four requirements hold for every sufficiently
large $Y$, so they can be met simultaneously.
Together with $\operatorname{Tr}A_{X_n}
=\operatorname{Tr}B_{X_n}=0$ and (MA-JP.17),
this already proves all four limits (MA-JP.15)
for sharp cutoffs.
It is an actual diagonal choice specified by four
trace-norm inequalities.  No unproved rate such as
the sufficiency of the exact power $Y=X^2$ is asserted.

\emph{The source's smooth symmetric cutoffs.}
Choose a fixed smooth $b:\mathbb R\to[0,1]$ with
$b(u)=0$ for $u\le0$ and $b(u)=1$ for $u\ge1$.
For fixed $X\ge2$ and sufficiently small $\delta>0$
put
\[
 \begin{aligned}
 \rho_{X,\delta}(t)
  &=b((X-t)/\delta)\,b((X-t^{-1})/\delta),\qquad t>0,\\
 b_{Y,\delta}(\xi)
  &=b((Y-|\xi|)/\delta).
 \end{aligned}
 \tag{MA-JP.21}
\]
The second function is smooth at zero because it
is identically one in a neighborhood of zero.
The first is smooth, multiplicatively symmetric,
supported on $[1/X,X]$, and equal to one on
$[1/X+\delta,X-\delta]$ when $X(X-\delta)\ge1$.
Indeed $1/(X-\delta)-1/X
=\delta/(X(X-\delta))\le\delta$.
Thus it is a cutoff of the exact source definition.
As $\delta\downarrow0$ both cutoffs converge boundedly
and pointwise almost everywhere to their sharp versions.

At the already fixed $X_n,Y_n$, the spatial trace-norm
convergences proved after (MA-JP.16) hold.
Moreover
$\mathcal F_e^{-1}M_{b_{Y_n,\delta}}\mathcal F_e$
converges strongly, with its adjoint, to $Q_{Y_n}$.
Splitting the difference of the two local products
at the fixed sharp $B_{X_n}$ and using (MA-JP.20)'s
trace-ideal argument shows that the entire operator
in the last line of (MA-JP.15) converges in trace norm to
$B_{X_n}D_SQ_{Y_n}D_S^*$ as $\delta\downarrow0$.
Every factor has its stated actual Hilbert domain.

Choose $\delta_n>0$ after $X_n,Y_n$ so that all
four resulting changes of trace have absolute
value at most $1/n$ and
\[
             \delta_n<
        \min\{1/(4X_n),\,1/(n(X_n+Y_n))\}.
 \tag{MA-JP.22}
\]
This proves (MA-JP.14)--(MA-JP.15), with
$\rho_n=\rho_{X_n,\delta_n}$ and
$b_n=b_{Y_n,\delta_n}$.
The smoothing does not change the band defining
the \emph{global} orthogonal PW projection; it
smooths the actual local additive Fourier cutoff.
Those are two distinct operations in the source
comparison, as their definitions require.

\emph{The literal Gram remainder.}
For the smoothed data put
$A_n=M_{\rho_n}C_g$ and use the closed
$T_{Y_n}$, Gram $G_{Y_n}$ and core-form operator
$K_{\rho_n,h}$ of MA-GP-008.
At each fixed $n$ that theorem gives
\[
 \begin{split}
 &\operatorname{Tr}_H
      \bigl(A_n(Q_{Y_n}^{\rm glob}-\Pi_{Y_n}^{S,0})\bigr)\\
 &\quad=\lim_{\eta\downarrow0}
    \operatorname{Tr}_{\mathcal P_{Y_n}^0}
     \left((G_{Y_n}+\eta)^{-1/2}
             K_{\rho_n,h}(G_{Y_n}+\eta)^{-1/2}
                    -(W_{Y_n}^S)^*A_nW_{Y_n}^S\right).
 \end{split}
 \tag{MA-JP.23}
\]
Only the bounded isometry $W_{Y_n}^S$ has changed
from (MA-GP.74), so the same cross-space cyclicity
proves this identity.  If a single fully finite
sequence is desired, choose $\eta_n>0$ at each
fixed $n$ so that the Gram trace differs from
its limit by at most $1/n$.
Then the trace of the fully bounded, trace-class
operator inside (MA-JP.23), with $\eta=\eta_n$,
tends to $-1$.  This retains both Gram resolvents
and is an actual nonvanishing finite-operator
remainder.
\end{proof}

\paragraph{What this path does and does not decide.}
The result disproves a rate-independent identification
of the two actual traces under only the joint
divergence and thin-smoothing conditions printed on
p.~519, when the annular repair is made.
It also prevents dropping the exact Gram correction
uniformly on that allowed collection of paths.
The path has $Y_n/X_n\to\infty$.
It does not refute the separately imposed local
repair $Y/X\to0$, nor prove a quantitative Gram
asymptotic on that different regime.
It does not exchange the two iterated limits of
MA-GP-009.
The nonzero quotient trace is selected from a
retained block, so this argument is different
from MA-EG's criterion involving discarded blocks.
In particular it is already instantiated by a
critical-line zero, independently of whether any
exterior zero exists.

The finite-prime set can meet the source's support
condition exactly: after $g$ is fixed, choose
$T>\max\{e^x,e^{-x}:x\in\operatorname{supp}g\}$ and
include every rational prime at most $T$ in $S_f$.
Then every omitted prime has norm greater than $T$.
The test is in the trivial compact isotype, so the
outside-unit invariance condition also holds.
The proof used only that this set is fixed and
finite; its actual $D_S$ was retained throughout.
The second line of (MA-JP.15) also retains the
actual global invariant PW family for that same
set $S$.  The explicit Gram-resolvent kernel
(MA-JP.23) remains the one-place parametrization
for which MA-GP-008 proves it; no uncomputed
mixed-place Gram kernel is assigned to the
second line.

\paragraph{Grothendieck placement and source scope.}
The witness is chosen in one finite-length
$\mathbb C[z]$-module of the actual completed
spectral quotient.  Its trace normalization is
the finite-dimensional trace with its retained
length.  The subsequent limits use actual
Hilbert trace ideals and the nuclear section
for the local quotient.  This is an instance of
the inherited cote~1 trace comparison
(scans~079--080 and~111--112, BR-TR-01/02);
it is not a new arithmetic positivity form,
an Artin-motive zero bound, or an identification
of the local and global Hilbert metrics.

The source-rate statement was checked in the
supplied Masullo text at Definition~9.3.3.1,
p.~427, the footnote to Theorem~12.4.1.1,
p.~514, (12.4.1.25)--(12.4.1.28), p.~519,
and the complete proof of Theorem~12.5.1.1,
p.~520.  These are fresh text readings;
no new rendered-page inspection is claimed here.
Hardy's primary theorem statement was checked
in its reproduced text at printed p.~1012.
Its full 1914 proof is used as a classical
published theorem, not presented as a newly
verified proof or numerical zero certificate.

Hardy, G.~H. (1914). Sur les z\'eros de la fonction
$\zeta(s)$ de Riemann. \emph{Comptes rendus
hebdomadaires des s\'eances de l'Acad\'emie
des sciences, 158}, 1012--1014.
Primary-page reproduction:
\url{https://fr.wikisource.org/wiki/Page:Comptes_rendus_hebdomadaires_des_s\%C3\%A9ances_de_l\%E2\%80\%99Acad\%C3\%A9mie_des_sciences,_tome_158,_1914.djvu/1014}.
